% problem_id: S001-spaces-morphisms-lemma-normalization
% source_id: S001 (Stacks Project)
% source_file: spaces-morphisms.tex
% statement_locator: spaces-morphisms.tex lines 10138--10167
% proof_locator: spaces-morphisms.tex lines 10169--10329
% env: lemma
% id_note: label 跨文件重名 ⇒ id 用文件限定形式
% pairing: tight (verified)
% status: complete (statement + author proof verbatim + reason + locators)
%
\begin{lemma}
\label{lemma-normalization}
Let $S$ be a scheme. For every algebraic space $X$ over $S$ satisfying
the equivalent conditions of Lemma \ref{lemma-prepare-normalization}
there exists a morphism of algebraic spaces
$$
\nu_X : X^\nu \longrightarrow X
$$
with the following properties
\begin{enumerate}
\item if $X$ satisfies the equivalent conditions of
Lemma \ref{lemma-prepare-normalization} then
$X^\nu$ is normal and $\nu_X$ is integral,
\item if $X$ is a scheme such that every quasi-compact open has finitely
many irreducible components, then $\nu_X : X^\nu \to X$ is the
normalization of $X$ constructed in
Morphisms, Section \ref{morphisms-section-normalization},
\item if $f : X \to Y$ is a morphism of algebraic spaces over $S$
which both satisfy the equivalent conditions of
Lemma \ref{lemma-prepare-normalization} and every codimension $0$
point of $X$ is mapped by $f$ to a codimension $0$ point of $Y$, then
there is a unique morphism $f^\nu : X^\nu \to Y^\nu$ of algebraic
spaces over $S$ such that $\nu_Y \circ f^\nu = f \circ \nu_X$, and
\item if $f : X \to Y$ is an \'etale or smooth morphism of algebraic
spaces and $Y$ satisfies the equivalent conditions of
Lemma \ref{lemma-prepare-normalization}, then the hypotheses of (3)
hold and the morphism $f^\nu$ induces an isomorphism
$X^\nu  \to X \times_Y Y^\nu$.
\end{enumerate}
\end{lemma}

\begin{proof}
Consider the category $\mathcal{C}$ whose objects are the
schemes $U$ over $S$ such that every quasi-compact open of
$U$ has finitely many irreducible components and whose morphisms
are those morphisms $g : U \to V$ of schemes over $S$ such
that every generic point of an irreducible component of $U$
is mapped to the generic point of an irreducible component of $V$.
We have already shown that
\begin{enumerate}
\item[(a)] for $U \in \Ob(\mathcal{C})$ we have a
normalization morphism $\nu_U : U^\nu \to U$ as in
Morphisms, Definition \ref{morphisms-definition-normalization},
\item[(b)] for $U \in \Ob(\mathcal{C})$ the morphism $\nu_U$
is integral and $U^\nu$ is a normal scheme, see
Morphisms, Lemma \ref{morphisms-lemma-normalization-normal},
\item[(c)] for every $g : U \to V \in \text{Arrows}(\mathcal{C})$
there is a unique morphism $g^\nu : U^\nu \to V^\nu$
such that $\nu_V \circ g^\nu = g \circ \nu_U$, see
Morphisms, Lemma \ref{morphisms-lemma-normalization-normal}
part (4) applied to the composition $X^\nu \to X \to Y$,
\item[(d)] if $V \in \Ob(\mathcal{C})$ and $g : U \to V$ is
\'etale or smooth, then $U \in \Ob(\mathcal{C})$ and
$g \in \text{Arrows}(\mathcal{C})$ and the morphism $g^\nu$
induces an isomorphism $U^\nu \to U \times_V V^\nu$, see
Lemma \ref{lemma-flat-locally-finite-type-points-codimension-0}
and More on Morphisms, Lemma
\ref{more-morphisms-lemma-normalization-and-smooth}.
\end{enumerate}
Our task is to extend this construction to the corresponding
category of algebraic spaces $X$ over $S$.

\medskip\noindent
Let $X$ be an algebraic space over $S$ satisfying
the equivalent conditions of Lemma \ref{lemma-prepare-normalization}.
Let $U \to X$ be a surjective \'etale morphism where $U$ is a scheme.
Set $R = U \times_X U$ with projections $s, t : R \to U$ and
$j = (t, s) : R \to U \times_S U$ so that $X = U/R$, see
Spaces, Lemma \ref{spaces-lemma-space-presentation}.
Observe that $U$ and $R$ are objects of $\mathcal{C}$
by our assumptions on $X$ and that the morphisms
$s$ and $t$ are \'etale morphisms of schemes over $S$.
By (a) we have the normalization morphisms
$\nu_U : U^\nu \to U$ and $\nu_R : R^\nu \to R$,
by (d) we have morphisms $s^\nu : R^\nu \to U^\nu$,
$t^\nu : R^\nu \to U^\nu$ which define
isomorphisms $R^\nu \to R \times_{s, U} U^\nu$ and
$R^\nu \to U^\nu \times_{U, t} R$. It follows that
$s^\nu$ and $t^\nu$ are \'etale (as they are isomorphic
to base changes of \'etale morphisms). The induced morphism
$j^\nu = (t^\nu, s^\nu) : R^\nu \to U^\nu \times_S U^\nu$
is a monomorphism as it is equal to the composition
\begin{align*}
R^\nu
& \to
(U^\nu \times_{U, t} R) \times_R (R \times_{s, U} U^\nu) \\
& =
U^\nu \times_{U, t} R \times_{s, U} U^\nu \\
& \xrightarrow{j}
U^\nu \times_U (U \times_S U) \times_U U^\nu \\
& =
U^\nu \times_S U^\nu
\end{align*}
The first arrow is the diagonal morphism of $\nu_R$.
(This tells us that $R^\nu$ is a subscheme of the restriction of
$R$ to $U^\nu$.) A formal computation with fibre products
using property (d) shows that
$R^\nu \times_{s^\nu, U^\nu, t^\nu} R^\nu$ is the normalization
of $R \times_{s, U, t} R$. Hence the \'etale morphism
$c : R \times_{s, U, t} R \to R$ extends uniquely to $c^\nu$ by (d).
The morphism $c^\nu$ is compatible with the projection
$\text{pr}_{13} : U^\nu \times_S U^\nu \times_S U^\nu \to U^\nu \times_S U^\nu$.
Similarly, there are morphisms $i^\nu : R^\nu \to R^\nu$
compatible with the morphism $U^\nu \times_S U^\nu \to U^\nu \times_S U^\nu$
which switches factors and there is a morphism $e^\nu : U^\nu \to R^\nu$
compatible with the diagonal morphism $U^\nu \to U^\nu \times_S U^\nu$.
All in all it follows that $j^\nu : R^\nu \to U^\nu \times_S U^\nu$
is an \'etale equivalence relation.
At this point we may and do set $X^\nu = U^\nu/R^\nu$
(Spaces, Theorem \ref{spaces-theorem-presentation}).
Then we see that we have $U^\nu = X^\nu \times_X U$ by
Groupoids, Lemma \ref{groupoids-lemma-criterion-fibre-product}.

\medskip\noindent
What have we shown in the previous paragraph is this: for every algebraic
space $X$ over $S$ satisfying the equivalent conditions of
Lemma \ref{lemma-prepare-normalization} if we choose a surjective \'etale
morphism $g : U \to X$ where $U$ is a scheme, then we obtain
a cartesian diagram
$$
\xymatrix{
X^\nu \ar[d]_{\nu_X} & U^\nu \ar[l]^{g^\nu} \ar[d]^{\nu_U} \\
X & U \ar[l]_g
}
$$
of algebraic spaces. This immediately implies that $X^\nu$ is a normal
algebraic space and that $\nu_X$ is an integral morphism. This gives
part (1) of the lemma.

\medskip\noindent
We will show below that the morphism $\nu_X : X^\nu \to X$
up to unique isomorphism is independent of the choice of $g$,
but for now, if $X$ is a scheme, we choose $\text{id} : X \to X$
so that it is clear that we have part (2) of the lemma.

\medskip\noindent
We still have to prove parts (3) and (4). Let $g : U \to X$ and
$\nu_X : X^\nu \to X$ and $g^\nu : U^\nu \to X^\nu$ be as above.
Let $Z$ be a normal scheme and let $h : Z \to U$ and $a : Z \to X^\nu$
be morphisms over $S$ such that $g \circ h = \nu_X \circ a$ and such that
every irreducible component of $Z$ dominates an irreducible component of $U$
(via $h$). By
Morphisms, Lemma \ref{morphisms-lemma-normalization-normal} part (4)
we obtain a unique morphism $h^\nu : Z \to U^\nu$ such
that $h = \nu_U \circ h^\nu$. Picture:
$$
\xymatrix{
X^\nu \ar[d]_{\nu_X} &
U^\nu \ar[l]^{g^\nu} \ar[d]^{\nu_U} &
Z \ar[l]^{h^\nu} \ar@/_1em/[ll]_a \ar[dl]^h
\\
X & U \ar[l]_g
}
$$
Observe that $a = g^\nu \circ h^\nu$. Namely, since the square with
corners $X^\nu$, $X$, $U^\nu$, $U$ is cartesian, this follows
immediately from the fact that $h^\nu$ is unique (given $h$).
In other words, given $h : Z \to U$ as above (and not $a$)
there is a unique morphism $a : Z \to X^\nu$ with
$\nu_X \circ a = g \circ h$.

\medskip\noindent
Let $f : X \to Y$ be as in part (3) of the statement of the lemma.
Suppose we have chosen surjective \'etale morphisms $U \to X$ and
$V \to Y$ where $U$ and $V$ are schemes such that $f$ lifts to a
morphism $g : U \to V$. Then $g \in \text{Arrows}(\mathcal{C})$
and we obtain a unique morphism $g^\nu : U^\nu \to V^\nu$ compatible
with $\nu_U$ and $\nu_V$. However, then the two morphisms
$$
R^\nu = U^\nu \times_{X^\nu} U^\nu
\to U^\nu \to V^\nu \to Y^\nu
$$
must be the same by our comments in the previous paragraph (applied
with $Y$ in stead of $X$).
Since $X^\nu$ is constructed by taking the quotient of
$U^\nu$ by $R^\nu$ it follows that we obtain a (unique)
morphism $f^\nu : X^\nu \to Y^\nu$ as stated in (3).

\medskip\noindent
To see that the construction of $X^\nu$ is independent of the choice
of $g : U \to X$ surjective \'etale, apply the construction in the
previous paragraph to $\text{id} : X \to X$ and a morphism
$U' \to U$ between \'etale coverings of $X$. This is enough because
given any two \'etale coverings of $X$ there is a third one which
dominates both. The reader shows that the morphism between the two
normalizations constructed using either $U' \to X$ or $U \to X$
becomes an isomorphism after base change to $U'$ and hence was
an isomorphism. We omit the details.

\medskip\noindent
We omit the proof of (4) which is similar; hint use part (d) above.
\end{proof}
